\chapter{实验与材料复现入口}

本附录给出正文图表所对应的源码、参数和公开数据入口。报告源文件位于
\nolinkurl{docs/reports/2026-10-22-reader-edit/}；同目录的
\nolinkurl{build-manifest.json}记录主文件、各章、图片和PDF的SHA-256摘要，
\nolinkurl{material-sources.json}记录所有生成表图的输入摘要。十种组合的实际运行源码提交为
\texttt{a202c6d3c87798911675b11a46b44903405370df}，源码归档摘要为
\texttt{acd3f84e655a29a7fac2b755889ff5c04ac128a995bb2664bf91d09c7ff614dc}。

\section{数据、参数与源码位置}

\begin{longtable}{@{}>{\raggedright\arraybackslash}p{.2\textwidth}>{\raggedright\arraybackslash}p{.76\textwidth}@{}}
\caption{实验数据与源码入口}\label{tab:reproduction-index}\\
\toprule
\rowcolor{black!8} 内容 & 仓库入口 \\
\midrule
十种组合同机比较 & \nolinkurl{docs/research/2026-10-21-ten-combination-benchmark-results/}；主表为\nolinkurl{basic-comparison/basic-comparison.csv}，公开逐项记录位于\nolinkurl{calibration-02/}和\nolinkurl{basic-supplement-01/} \\
A/B多根及通信量 & \nolinkurl{docs/deliveries/2026-10-19-report-experiment-supplement/}，其中\nolinkurl{master-table.json}汇总时间和对象尺寸，\nolinkurl{evidence/benchmark-A.json}与\nolinkurl{evidence/benchmark-B.json}保存逐根观测 \\
Sageloli结构化采样 & \nolinkurl{docs/research/team-b-delivery/2026-10-19-sageloli-ring-target/}；两档运行记录位于\nolinkurl{evidence/param-I-fresh.json}与\nolinkurl{evidence/param-II-fresh.json} \\
SRGG紧凑编码 & \nolinkurl{docs/research/team-b-delivery/2026-10-18-srgg-compact-benchmark/}；对象结果为\nolinkurl{wire.csv}，协议载荷功能记录位于\nolinkurl{docs/research/2026-10-19-srgg-paper-latte/} \\
Lai接收会话 & \nolinkurl{docs/research/2026-10-20-multi-backend-optimization/verified-summary.json}；C组条目记录原路径、准备和会话内解密中位时间 \\
Wang--PW目标档 & \nolinkurl{docs/research/2026-10-19-wang-pw-target-s0-4/}；公开说明、分阶段结果和验证入口见\nolinkurl{PUBLICATION_zh-CN.md}及\nolinkurl{publication/} \\
完整参数配置 & 三组代表位于\nolinkurl{docs/deliveries/2026-10-19-report-three-selection/suites/}；十种组合矩阵位于\nolinkurl{docs/research/2026-10-21-ten-combination-benchmark-plan/matrix.json} \\
\bottomrule
\end{longtable}

\section{报告图表与PDF重建}

在仓库根目录运行以下命令。Python环境需要NumPy、Matplotlib、pypdf和pypdfium2；
PDF编译使用Tectonic、XeLaTeX或latexmk。中文图表字体通过\texttt{--font}指定。

\subsection{根据归档数据重建图表与报告}
\begingroup\footnotesize
\begin{verbatim}
python docs/reports/2026-10-22-reader-edit/prepare_materials.py
python docs/reports/2026-10-22-reader-edit/build_pdf.py
python docs/reports/2026-10-22-reader-edit/check_report.py --require-pdf
\end{verbatim}
\endgroup

第一条命令从JSON与CSV重建数据表、数据图和来源摘要；第二条调用TeX引擎生成PDF与输入摘要；
第三条检查文档完整性、引文、图表数据和编译日志。视觉复核记录位于
\nolinkurl{visual-review.json}。字体与编译器路径可分别通过\texttt{--font}和\texttt{--engine}传入。

十种组合的归档结果由比较工具重新汇总；以下为Bash命令：
\begingroup\footnotesize
\begin{verbatim}
python -m experiments.research.ten_combination_basic_comparison \
  --design docs/research/2026-10-21-ten-combination-benchmark-results/\
basic-design.json \
  --reused-public docs/research/2026-10-21-ten-combination-benchmark-results/\
calibration-02 \
  --supplemental-public docs/research/2026-10-21-ten-combination-benchmark-results/\
basic-supplement-01 --output <新输出目录> --plots
\end{verbatim}
\endgroup

\subsection{重新运行密码实验}
实际运行入口为\nolinkurl{experiments/research/ten_combination_benchmark.py}及其工作进程。
运行时记录源码提交、原生组件、依赖、线程和硬件信息，采用新的私有状态目录保存密钥。
详细配置和运行步骤见对应实验目录的README。正文时间数据对应表~\ref{tab:reproduction-index}所列归档批次。
